\documentclass[journal]{IEEEtran}
\usepackage[a5paper, margin=10mm, onecolumn]{geometry}
\usepackage[cmex10]{amsmath}
\usepackage{amssymb,amsfonts,amsthm}
\usepackage{gvv-book}
\usepackage{gvv}
\usepackage{hyperref}

\begin{document}


\title{2.10.54}
\author{EE25BTECH11025 - Ganachari Vishwambhar}
\maketitle

\textbf{Question}:\newline
let $\vec{a, b,c}$ be unit vectors such that $\vec{a}+\vec{b}+\vec{c}=\vec{0}$. Which of the following are correct?\\
\begin{enumerate}
    \item $\vec{a}\times \vec{b}=\vec{b}\times \vec{c}=\vec{c}\times \vec{a}=\vec{0}$
    \item $\vec{a}\times \vec{b}=\vec{b}\times \vec{c}=\vec{c}\times \vec{a}\neq \vec{0}$
    \item $\vec{a}\times \vec{b}=\vec{b}\times \vec{c}=\vec{a}\times \vec{c}\neq \vec{0}$
    \item $\vec{a}\times \vec{b},\vec{b}\times \vec{c}, \vec{c}\times \vec{a}$ are mutually perpendicular.
\end{enumerate}
\textbf{Solution: }\\

Given:
\begin{align}
    \vec{a}+\vec{b}+\vec{c}=0\\
    \vec{c}=-\brak{\vec{a}+\vec{b}}\\
    \vec{b}=-\brak{\vec{a}+\vec{c}}\\
    \vec{a}=-\brak{\vec{b}+\vec{c}}
\end{align}

Since $\vec{c}$ is a unit vector:
\begin{align}
    ||\vec{c}||=1\\
\end{align}

From(2):
\begin{align}
    1=||\vec{c}||^2=\brak{\vec{a}+\vec{b}}^\top\brak{\vec{a}+\vec{b}}\\
    1=\vec{a}^\top \vec{a}+\vec{b}^\top \vec{b}+2\vec{a}^\top \vec{b}\\
    1=1+1+2\vec{a}^\top \vec{b}\\
    \vec{a}^\top \vec{b}=\frac{-1}{2}
\end{align}

Similarly,
\begin{align}
    \vec{b}^\top \vec{c}=\frac{-1}{2}\\
    \vec{c}^\top \vec{a}=\frac{-1}{2}
\end{align}

$\therefore$Angle between the vectors is neither $0\degree$ nor $180\degree$:
\begin{align}
    \vec{a}\times \vec{b}\neq 0\\
    \vec{b}\times \vec{c}\neq 0\\
    \vec{c}\times \vec{a}\neq 0
\end{align}

Proving $\vec{b}\times \vec{c}=\vec{a}\times \vec{b}$.\\

From(2):
\begin{align}
    \vec{b}\times \vec{c}=\vec{b}\times \brak{-\brak{\vec{a}+\vec{b}}}\\
    \vec{b}\times \vec{c}=-\brak{\vec{b}\times \vec{a}}-\brak{\vec{b}\times \vec{b}}\\
    \vec{b}\times \vec{c}=\vec{a}\times \vec{b}\\
    \because \vec{b}\times \vec{b}=0
\end{align}

Similarly,
\begin{align}
    \vec{a}\times \vec{b}=\vec{b}\times \vec{c}.\\
    \vec{c}\times \vec{a}=\vec{a}\times \vec{b}.
\end{align}

From (13, 14, 15, 18, 20, 21), we can conclude that:
\begin{align}
    \vec{a}\times \vec{b}=\vec{b}\times \vec{c}=\vec{c}\times \vec{a}\neq 0
\end{align}

Hence, option (3) is correct.

\begin{figure}[h!]
   \centering
   \includegraphics[width=0.7\linewidth]{figs/plot.png}
   \caption{Plot of the given line, point and reflected point}
   \label{}
\end{figure}
\end{document}  